% Clase 20 --- Corrientes parásitas: el péndulo de Foucault y los núcleos
% laminados (§3).
% "Péndulo de bronce macizo que se frena entre bobinas encendidas; con cortes
%  (peine), oscila libre."
% Dos paneles con la misma placa, el mismo imán y la misma velocidad: lo único
% que cambia son los cortes. Es el patrón de "una sola variable en discusión".
\begin{center}
\begin{tikzpicture}[scale=1]

  \providecommand{\montaje}{%
    \draw[figgray, line width=0.9pt] (0,3.2) -- (0,3.4);
    \node[figgray, circle, fill=figgray, inner sep=1.3pt] at (0,3.4) {};
    \draw[figgray, line width=0.9pt] (0,3.4) -- (1.15,0.9);
    \fill[figblue!8] (-1.5,-1.3) rectangle (1.5,0.5);
    \draw[figgray, line width=0.7pt] (-1.5,-1.3) rectangle (1.5,0.5);
    \foreach \xx in {-1.1,-0.4,0.3,1.0} {
      \foreach \yy in {-1.0,-0.4,0.2} {
        \node[figblue, scale=0.55] at (\xx,\yy) {$\otimes$};
      }}
    \node[figlblaux, anchor=south east] at (1.5,0.55) {$\vec B$};}

  % =============== (a) placa maciza ===============
  \begin{scope}
    \montaje
    \fill[figamber!30] (0.55,-0.75) rectangle (1.75,0.9);
    \draw[figamber, line width=1.1pt] (0.55,-0.75) rectangle (1.75,0.9);
    % corrientes parásitas
    \draw[figvecres, line width=0.9pt] (1.15,0.05) circle (0.32);
    \draw[figvecres, line width=0.9pt] (1.15,0.05) circle (0.55);
    \node[figlbl, figred, anchor=west] at (1.85,0.05) {corrientes};
    \draw[figvecres, line width=1.3pt] (0.35,-1.05) -- (-0.55,-1.05);
    \node[figlbl, figred, anchor=north] at (-0.1,-1.15) {frena};
    \node[figlbl, anchor=north, align=center] at (0,-1.85)
      {(a) placa maciza:\\[-0.2em] se detiene enseguida};
  \end{scope}

  % =============== (b) placa con cortes ===============
  \begin{scope}[xshift=6.6cm]
    \montaje
    \fill[figamber!30] (0.55,-0.75) rectangle (1.75,0.9);
    \draw[figamber, line width=1.1pt] (0.55,-0.75) rectangle (1.75,0.9);
    \foreach \xx in {0.79,1.03,1.27,1.51} {
      \draw[white, line width=2.2pt] (\xx,-0.75) -- (\xx,0.62);
      \draw[figamber, line width=0.5pt] (\xx,-0.75) -- (\xx,0.62);
    }
    \node[figlblaux, figred, anchor=west, align=left] at (1.85,0.05)
      {sin lugar para\\[-0.25em] los remolinos};
    \node[figlbl, anchor=north, align=center] at (0,-1.85)
      {(b) placa ranurada:\\[-0.2em] oscila casi libre};
  \end{scope}

\end{tikzpicture}\\[0.5em]
{\small\itshape Al entrar y salir de la región con campo, el flujo a través de
la placa cambia y se inducen corrientes en el propio metal ---corrientes
\textbf{parásitas} o de Foucault, remolinos cerrados dentro del material---.
Lenz vuelve a decidir el signo: se oponen al movimiento, de modo que la placa
maciza se frena en pocas oscilaciones y la energía mecánica termina como calor
por efecto Joule. La versión ranurada es el control experimental: es el mismo
metal, la misma masa y el mismo campo, pero los cortes interrumpen los caminos
cerrados y los remolinos no pueden formarse, así que el péndulo oscila casi sin
pérdidas. De acá sale, sin ningún paso adicional, la razón por la que los
núcleos de los transformadores se construyen \textbf{laminados} y no macizos:
las chapas aisladas entre sí cumplen el papel de las ranuras, cortan las
corrientes parásitas y evitan que buena parte de la energía se disipe en
calentar el hierro. Y del lado en que sí se las quiere, son las mismas
corrientes que frenan un tren sin tocarlo.}
\end{center}
